\chapter{Applications linéaires}\label{ch:b1:linmaps}

Entre \omterm{def:b1:vspaces:def}{espaces vectoriels}, les \omterm{def:b1:logic:map}{applications} qui méritent d'être
étudiées sont celles qui respectent la structure~: les
\emph{\omterm{def:b1:linmaps:def}{applications linéaires}}. Leurs deux \omterm{def:b1:vspaces:subspace}{sous-espaces} fondamentaux
--- \omterm{def:b1:linmaps:kerim}{noyau} et \omterm{def:b1:linmaps:kerim}{image} --- mesurent l'\omterm{def:b1:logic:inj}{injectivité} et la \omterm{def:b1:logic:inj}{surjectivité}, et
en \omterm{def:b1:findim:def}{dimension finie} le \omterm{thm:b1:linmaps:ranknullity}{théorème du rang} relie leurs tailles par une
loi de conservation. \omterm{def:b1:linmaps:projection}{Projections} et \omterm{def:b1:linmaps:projection}{symétries}, puis \omterm{def:b1:linmaps:forms}{formes linéaires}
et \omterm{def:b1:linmaps:forms}{hyperplans}, achèvent le chapitre.

\section{Définitions et premières propriétés}

\begin{definition}\label{def:b1:linmaps:def}
Soient $E, F$ des $K$-espaces vectoriels. Une \omterm{def:b1:logic:map}{application} $u \colon E
\to F$ est \emph{linéaire}\index{application linéaire} lorsque
\[
\forall x, y \in E,\ \forall \lambda \in K, \qquad
u(x + \lambda y) = u(x) + \lambda u(y).
\]
On a alors $u(0) = 0$ et $u(\sum \lambda_i x_i) = \sum \lambda_i u(x_i)$.
L'\omterm{def:b1:logic:sets}{ensemble} $\mathcal{L}(E, F)$ des \omterm{def:b1:logic:map}{applications} linéaires est lui-même
un \omterm{def:b1:vspaces:def}{espace vectoriel}~; la composée d'\omterm{def:b1:logic:map}{applications} linéaires est
linéaire, et bilinéaire en le couple. Un \emph{endomorphisme} est une
\omterm{def:b1:logic:map}{application} linéaire $E \to E$~; un \emph{isomorphisme} est une
\omterm{def:b1:logic:map}{application} linéaire \omterm{def:b1:logic:inj}{bijective} (sa réciproque est alors automatiquement
linéaire)~; l'\omterm{def:b1:logic:map}{application} $u^{-1}$ d'un isomorphisme, les composées
d'isomorphismes, sont des isomorphismes.
\end{definition}

\begin{proof}[Démonstration du fait que la réciproque est linéaire]
Soient $u$ \omterm{def:b1:linmaps:def}{linéaire} \omterm{def:b1:logic:inj}{bijective}, $y, y' \in F$ et $\lambda \in K$.
Posons $x = u^{-1}(y)$ et $x' = u^{-1}(y')$. Alors
\[
u\bigl(x + \lambda x'\bigr) = u(x) + \lambda u(x')
= y + \lambda y' ,
\]
et en appliquant $u^{-1}$ aux deux extrémités~: $u^{-1}(y + \lambda y') = x
+ \lambda x' = u^{-1}(y) + \lambda\,u^{-1}(y')$. Rien de
$u^{-1}$ n'a été calculé~: la linéarité se transporte par la seule
propriété caractéristique de $u$ --- un schéma à retenir,
car la structure voyage souvent gratuitement le long des
\omterm{def:b1:logic:inj}{bijections}.
\end{proof}

\begin{proposition}[Une application linéaire est connue sur une base]\label{prop:b1:linmaps:basis}
Soient $(e_1, \dots, e_n)$ une \omterm{def:b1:vspaces:free}{base} de $E$ et $(v_1, \dots, v_n)$
des vecteurs quelconques de $F$. Il existe exactement une \omterm{def:b1:linmaps:def}{application
linéaire} $u \colon E \to F$ telle que $u(e_i) = v_i$ pour tout $i$. De
plus~:
\[
u \text{ injective} \iff (v_i) \text{ libre};
\qquad
u \text{ surjective} \iff (v_i) \text{ engendre } F .
\]
\end{proposition}

\begin{proof}
Existence et unicité~: tout $x$ a des \omterm{prop:b1:vspaces:coordinates}{coordonnées} uniques $x = \sum
\lambda_i e_i$ (\cref{prop:b1:vspaces:coordinates})~; la linéarité
impose $u(x) = \sum \lambda_i v_i$, et cette formule définit bien une
\omterm{def:b1:linmaps:def}{application linéaire}.

\omterm{def:b1:logic:inj}{Injectivité}~: d'après la \cref{prop:b1:linmaps:kernel} ci-dessous, $u$ est
\omterm{def:b1:logic:inj}{injective} si et seulement si son \omterm{def:b1:linmaps:kerim}{noyau} est réduit à $\{0\}$. Or $u\bigl(\sum\lambda_i
e_i\bigr) = 0$ signifie exactement $\sum\lambda_i v_i = 0$. Si $(v_i)$
est \omterm{def:b1:vspaces:free}{libre}, cela force tous les $\lambda_i = 0$, c'est-à-dire que le \omterm{def:b1:linmaps:kerim}{noyau}
se réduit à $0$~: $u$ est \omterm{def:b1:logic:inj}{injective}. Si $(v_i)$ est liée, une relation
non triviale $\sum\lambda_i v_i = 0$ produit le vecteur non nul
$\sum\lambda_i e_i$ dans le \omterm{def:b1:linmaps:kerim}{noyau} ($(e_i)$ est \omterm{def:b1:vspaces:free}{libre})~: $u$ n'est pas
\omterm{def:b1:logic:inj}{injective}. Les deux conditions se correspondent terme à terme.

\omterm{def:b1:logic:inj}{Surjectivité}~: l'\omterm{def:b1:linmaps:kerim}{image} de $u$ est l'\omterm{def:b1:logic:sets}{ensemble} des $\sum\lambda_i
v_i$, c'est-à-dire exactement $\operatorname{Vect}(v_1, \dots, v_n)$, qui
vaut $F$ si et seulement si la famille est génératrice.
\end{proof}

\begin{definition}[Noyau et image]\label{def:b1:linmaps:kerim}
Pour $u \in \mathcal{L}(E, F)$~:
\[
\ker u = \{x \in E : u(x) = 0\} \subseteq E,
\qquad
\operatorname{im} u = u(E) \subseteq F ,
\]
tous deux \omterm{def:b1:vspaces:subspace}{sous-espaces vectoriels} (vérification directe avec le
critère).\index{noyau}\index{image}
\end{definition}

\begin{method}[Noyau et image, en pratique]
\label{met:b1:linmaps:kerim}
\emph{\omterm{def:b1:linmaps:kerim}{Noyau}}~: écrire $u(x) = 0$ comme un système portant sur les
\omterm{prop:b1:vspaces:coordinates}{coordonnées} (ou les coefficients) de $x$, le résoudre et le
paramétrer --- le \omterm{def:b1:linmaps:kerim}{noyau} sort avec une \omterm{def:b1:vspaces:free}{base} attachée
(\cref{met:b1:findim:computedim}). \emph{\omterm{def:b1:linmaps:kerim}{Image}}~: c'est le
\omterm{def:b1:vspaces:span}{sous-espace engendré} par les \omterm{def:b1:linmaps:kerim}{images} de \emph{n'importe quelle} famille
génératrice de $E$ --- en général une \omterm{def:b1:vspaces:free}{base}, donc $\operatorname{im} u =
\operatorname{Vect}\bigl(u(e_1), \dots, u(e_n)\bigr)$~; on élimine
ensuite les \omterm{def:b1:linmaps:kerim}{images} redondantes pour en extraire une \omterm{def:b1:vspaces:free}{base}. \emph{Raccourci}~:
calculer celui des deux qui est le plus facile et obtenir gratuitement la
dimension de l'autre par le \omterm{thm:b1:linmaps:ranknullity}{théorème du rang}
(\cref{thm:b1:linmaps:ranknullity})~; lorsqu'on connaît un
\omterm{def:b1:vspaces:subspace}{sous-espace} candidat plausible pour l'\omterm{def:b1:linmaps:kerim}{image}, la comparaison des dimensions
transforme l'inclusion facile en égalité
(\cref{thm:b1:findim:subspaces}). Les deux raccourcis sont utilisés dans
l'\cref{ex:b1:linmaps:delta} ci-dessous.
\end{method}

\begin{proposition}\label{prop:b1:linmaps:kernel}
$u$ est \omterm{def:b1:logic:inj}{injective} $\iff$ $\ker u = \{0\}$~; $u$ est \omterm{def:b1:logic:inj}{surjective} $\iff$
$\operatorname{im} u = F$.
\end{proposition}

\begin{proof}
Comme pour les \omterm{def:b1:structures:group}{groupes} (\cref{prop:b1:structures:kernel})~: $u(x) = u(y) \iff
u(x - y) = 0 \iff x - y \in \ker u$. Le second point est la
définition.
\end{proof}

\section{Le théorème du rang}

\begin{definition}\label{def:b1:linmaps:rank}
Le \emph{rang}\index{rang!d'une application linéaire} de $u \in \mathcal{L}(E,
F)$ (avec $E$ \omterm{def:b1:findim:def}{de dimension finie}) est $\operatorname{rk} u = \dim
\operatorname{im} u$ --- c'est aussi le rang de la famille
$\bigl(u(e_1), \dots, u(e_n)\bigr)$ pour n'importe quelle \omterm{def:b1:vspaces:free}{base} $(e_i)$ de $E$.
\end{definition}

\begin{theorem}[Théorème du rang]\label{thm:b1:linmaps:ranknullity}
Soient $E$ \omterm{def:b1:findim:def}{de dimension finie} et $u \in \mathcal{L}(E, F)$.
Alors\index{théorème du rang}
\[
\dim E = \dim \ker u + \operatorname{rk} u .
\]
Plus précisément, si $S$ est un \omterm{def:b1:vspaces:sum}{supplémentaire} quelconque de $\ker u$ dans
$E$, alors $u$ induit par restriction un \emph{isomorphisme} de $S$ sur
$\operatorname{im} u$.
\end{theorem}

\begin{proof}
Soit $S$ tel que $E = \ker u \oplus S$
(\cref{thm:b1:findim:subspaces}), et soit $v \colon S \to
\operatorname{im} u$ la restriction de $u$.

$v$ est \omterm{def:b1:logic:inj}{injective}~: $\ker v = S \cap \ker u = \{0\}$.

$v$ est \omterm{def:b1:logic:inj}{surjective}~: tout $u(x)$ avec $x = k + s$ ($k \in \ker u$, $s
\in S$) vaut $u(s) = v(s)$.

Donc $v$ est un isomorphisme~; un isomorphisme envoie une \omterm{def:b1:vspaces:free}{base} sur une \omterm{def:b1:vspaces:free}{base}
(\cref{prop:b1:linmaps:basis}), donc $\dim S = \dim\operatorname{im}
u$, et $\dim E = \dim\ker u + \dim S$ conclut.
\end{proof}

\begin{example}[Construire une application sur cahier des charges]
\label{ex:b1:linmaps:tospec}
Construisons $u \in \mathcal{L}(\R^3)$ avec $\ker u =
\operatorname{Vect}(1,1,1)$ et $\operatorname{im} u = \{z =
0\}$. Vérification d'abord~: le \omterm{thm:b1:linmaps:ranknullity}{théorème du rang} exige $1 + 2 = 3$ ---
c'est cohérent, donc une solution peut exister. Choisissons une \omterm{def:b1:vspaces:free}{base} adaptée
au \omterm{def:b1:linmaps:kerim}{noyau}, disons $\bigl((1,1,1),\ e_1,\ e_2\bigr)$
(\cref{ex:b1:findim:completion}), et prescrivons les \omterm{def:b1:linmaps:kerim}{images}
(\cref{prop:b1:linmaps:basis})~:
\[
u(1,1,1) = 0, \qquad u(e_1) = e_1, \qquad u(e_2) = e_2 .
\]
Alors $\ker u \supseteq \operatorname{Vect}(1,1,1)$ et
$\operatorname{im} u = \operatorname{Vect}(e_1, e_2) = \{z =
0\}$~; le \omterm{thm:b1:linmaps:ranknullity}{théorème du rang} force $\dim\ker u = 1$, donc le \omterm{def:b1:linmaps:kerim}{noyau} est
exactement la droite prescrite. Explicitement, en décomposant $(x, y, z)
= z(1,1,1) + (x - z)e_1 + (y - z)e_2$~:
\[
u(x, y, z) = (x - z,\ y - z,\ 0).
\]
La recette se généralise~: une \omterm{def:b1:linmaps:def}{application linéaire} de \omterm{def:b1:linmaps:kerim}{noyau} prescrit $N$
et d'\omterm{def:b1:linmaps:kerim}{image} prescrite $I$ existe exactement lorsque $\dim N + \dim I = \dim E$
--- la nécessité est le \omterm{thm:b1:linmaps:ranknullity}{théorème du rang}, la suffisance est cette
construction.
\end{example}

\begin{corollary}\label{cor:b1:linmaps:samedim}
Si $\dim E = \dim F$ (finie), alors pour $u \in \mathcal{L}(E, F)$~:
\[
u \text{ injective} \iff u \text{ surjective} \iff u \text{
bijective}.
\]
En particulier cela vaut pour les endomorphismes en \omterm{def:b1:findim:def}{dimension finie}. (C'est
faux en dimension infinie~: sur $K[X]$, la dérivation est \omterm{def:b1:logic:inj}{surjective}
mais non \omterm{def:b1:logic:inj}{injective}, et $P \mapsto XP$ est \omterm{def:b1:logic:inj}{injective} mais non
\omterm{def:b1:logic:inj}{surjective}.)
\end{corollary}

\begin{proof}
\omterm{def:b1:logic:inj}{Injective} $\iff \dim\ker u = 0 \iff \operatorname{rk} u = \dim E =
\dim F \iff \operatorname{im} u = F$ (un \omterm{def:b1:vspaces:subspace}{sous-espace} de dimension égale à
celle de l'espace est l'espace tout entier, \cref{thm:b1:findim:subspaces})
$\iff$ \omterm{def:b1:logic:inj}{surjective}.
\end{proof}

\begin{example}[L'interpolation, structurellement]\label{ex:b1:linmaps:interpolation}
Fixons $x_0, \dots, x_n$ distincts et soit $u \colon \R_n[X] \to
\R^{n+1}$, $P \mapsto (P(x_0), \dots, P(x_n))$~: elle est \omterm{def:b1:linmaps:def}{linéaire}. Son \omterm{def:b1:linmaps:kerim}{noyau}
est $\{P : \deg P \leq n,\ n+1 \text{ racines}\} = \{0\}$
(\cref{cor:b1:poly:nroots}). Dimensions égales $n + 1$~: $u$ est un
isomorphisme --- existence \emph{et} unicité du \omterm{def:b1:poly:def}{polynôme}
d'\omterm{thm:b1:poly:lagrange}{interpolation de Lagrange} (\cref{thm:b1:poly:lagrange}) en une ligne.

Le même schéma en une ligne traite des données mêlant valeurs et
\omterm{def:b1:derivative:def}{dérivées}~: $v \colon \R_3[X] \to \R^4$, $P \mapsto \bigl(P(0),
P'(0), P(1), P'(1)\bigr)$ est \omterm{def:b1:linmaps:def}{linéaire}, et son \omterm{def:b1:linmaps:kerim}{noyau} est formé des
\omterm{def:b1:poly:def}{polynômes} de degré $\leq 3$ ayant une racine double en $0$ \emph{et} en
$1$, c'est-à-dire divisibles par $X^2(X-1)^2$, de degré $4$~: seul $P =
0$ convient. Dimensions égales à nouveau~: tout quadruplet de données $(P(0),
P'(0), P(1), P'(1))$ est réalisé par exactement une cubique ---
l'interpolation d'Hermite, acquise par un calcul de \omterm{def:b1:linmaps:kerim}{noyau} avant
qu'aucune formule ne soit écrite (le devoir maison du \cref{ch:b1:det}
y rencontre son déterminant).
\end{example}

\begin{example}[Le théorème du rang à l'œuvre~: l'opérateur de différence]
\label{ex:b1:linmaps:delta}
Soit $\Delta \colon \R_n[X] \to \R_n[X]$, $P \mapsto P(X+1) -
P(X)$~: elle est \omterm{def:b1:linmaps:def}{linéaire}. \omterm{def:b1:linmaps:kerim}{Noyau}~: si $\Delta P = 0$, alors $P(0) = P(1) =
P(2) = \dots$, donc $P - P(0)$ a une infinité de racines et est
nul (\cref{cor:b1:poly:nroots})~: $\ker\Delta$ est la droite des
constantes. \omterm{thm:b1:linmaps:ranknullity}{Théorème du rang}~: $\operatorname{rk}\Delta = (n + 1) - 1
= n$. Comme $\deg \Delta P < \deg P$ pour $P$ non constant (les termes
de plus haut degré s'annulent), $\operatorname{im}\Delta \subseteq \R_{n-1}[X]$,
qui est de dimension exactement $n$~: l'inclusion est une égalité.
Conclusion, sans aucun calcul d'\omterm{def:b1:logic:map}{images réciproques}~: \emph{tout}
\omterm{def:b1:poly:def}{polynôme} $Q$ de degré $\leq n - 1$ est une différence $Q = P(X+1)
- P(X)$ --- la primitive discrète existe. (Comparer avec le
devoir maison du \cref{ch:b1:vspaces}, où $\Delta$ était
inversé explicitement dans la \omterm{def:b1:vspaces:free}{base} binomiale.)
\end{example}

\begin{example}[Comptabilité du rang le long d'une composée]
\label{ex:b1:linmaps:rankbookkeeping}
Sur $\R_2[X]$, composons la dérivation $D(P) = P'$ (de \omterm{def:b1:linmaps:rank}{rang} $2$~:
\omterm{def:b1:linmaps:kerim}{image} $\R_1[X]$, \omterm{def:b1:linmaps:kerim}{noyau} les constantes) avec elle-même. Alors $D
\circ D = D^2$ envoie $P \mapsto P''$, d'\omterm{def:b1:linmaps:kerim}{image} $\R_0[X]$~: de \omterm{def:b1:linmaps:rank}{rang}
$1$. Comparons avec les bornes générales~: la borne grossière donne
$\operatorname{rk} D^2 \leq \min(2, 2) = 2$~; la formule exacte
de l'\cref{exo:b1:linmaps:12} rend compte de la perte avec précision,
\[
\operatorname{rk} D^2 = \operatorname{rk} D -
\dim\bigl(\ker D \cap \operatorname{im} D\bigr) = 2 - 1 = 1 ,
\]
puisque les constantes (\omterm{def:b1:linmaps:kerim}{noyau} du $D$ extérieur) sont contenues dans
$\R_1[X]$ (\omterm{def:b1:linmaps:kerim}{image} du $D$ \omterm{def:b1:topology:closure}{intérieur}) avec la dimension $1$. Le \omterm{def:b1:linmaps:rank}{rang} se
perd exactement là où le \omterm{def:b1:linmaps:kerim}{noyau} extérieur tend une embuscade à l'\omterm{def:b1:linmaps:kerim}{image}
intérieure --- la phrase à retenir quand les \omterm{def:b1:linmaps:rank}{rangs} des composées se
comportent mal.
\end{example}

\begin{example}[Noyau et image de l'opérateur d'Euler]
\label{ex:b1:linmaps:euler}
Sur $\R_n[X]$, soit $u(P) = X\,P'$ (\omterm{def:b1:linmaps:def}{linéaire}~: la dérivation et la
multiplication par $X$ le sont). \emph{\omterm{def:b1:linmaps:kerim}{Noyau}}~: $XP' = 0$ force $P' =
0$ (un produit de \omterm{def:b1:poly:def}{polynômes} est nul seulement si l'un des facteurs l'est),
donc $\ker u$ est la droite des constantes. \emph{\omterm{def:b1:linmaps:kerim}{Image}}~: sur la
\omterm{def:b1:vspaces:free}{base} monomiale,
\[
u(X^k) = k\,X^{k} \qquad (k = 0, 1, \dots, n),
\]
donc $\operatorname{im} u = \operatorname{Vect}(X, 2X^2, \dots,
nX^n) = \operatorname{Vect}(X, X^2, \dots, X^n)$~: les
\omterm{def:b1:poly:def}{polynômes} de terme constant nul. Vérification avec le \omterm{thm:b1:linmaps:ranknullity}{théorème du rang}~:
$\operatorname{rk} u = (n + 1) - 1 = n$, ce qui est bien la
dimension trouvée. Deux remarques méritent d'être retenues. D'abord, ici
$\operatorname{im} u \oplus \ker u = \R_n[X]$ --- mais c'est un
\emph{heureux hasard} propre à cet opérateur, pas un théorème~: pour
$v(P) = P'$ sur $\R_1[X]$, qui ressemble à un décalage, $\ker v = \operatorname{im}
v = \R_0[X]$ et la somme n'est pas directe. Ensuite, la relation
$u(X^k) = kX^k$ dit que chaque monôme est simplement multiplié par un
scalaire sous $u$ --- une \omterm{def:b1:vspaces:free}{base} adaptée à l'\omterm{def:b1:logic:map}{application}, germe de l'idée
de valeur propre développée dans le volume de Licence~2.
\end{example}

\section{Projections et symétries}

\begin{definition}\label{def:b1:linmaps:projection}
Soit $E = F \oplus G$. La \emph{projection sur $F$ parallèlement à
$G$}\index{projection} envoie $x = f + g$ (décomposition unique) sur
$p(x) = f$~; la \emph{symétrie}\index{symétrie} associée est $s(x)
= f - g$. Toutes deux sont \omterm{def:b1:linmaps:def}{linéaires}, et $s = 2p - \mathrm{id}$.
\end{definition}

\begin{theorem}[Caractérisation algébrique]\label{thm:b1:linmaps:projchar}
\begin{enumerate}
  \item Un endomorphisme $p$ est une \omterm{def:b1:linmaps:projection}{projection} (sur un certain $F$
        parallèlement à un certain $G$) si et seulement si $p \circ p = p$~; alors $F =
        \operatorname{im} p = \ker(p - \mathrm{id})$ et $G = \ker
        p$.
  \item Un endomorphisme $s$ est une \omterm{def:b1:linmaps:projection}{symétrie} si et seulement si $s \circ s
        = \mathrm{id}$~; alors $E = \ker(s - \mathrm{id}) \oplus
        \ker(s + \mathrm{id})$.
\end{enumerate}
\end{theorem}

\begin{proof}
(1) Une \omterm{def:b1:linmaps:projection}{projection} vérifie $p(f + g) = f$ et $p(f) = f$~: $p^2 = p$.
Réciproquement, soit $p^2 = p$~; posons $F = \operatorname{im} p$, $G = \ker
p$. Tout $x$ s'écrit $x = p(x) + (x - p(x))$ avec $p(x) \in F$ et
$p\bigl(x - p(x)\bigr) = p(x) - p^2(x) = 0$~: $E = F + G$. Si $y \in
F \cap G$~: $y = p(z)$ et $p(y) = 0$, donc $y = p(z) = p^2(z) = p(y) =
0$~: la somme est directe, et $p$ est la \omterm{def:b1:linmaps:projection}{projection} sur $F$ parallèlement à $G$.
Enfin sur $F$~: $y = p(z)$ donne $p(y) = y$, donc $F \subseteq
\ker(p - \mathrm{id})$, et réciproquement $p(y) = y$ place $y$ dans
l'\omterm{def:b1:linmaps:kerim}{image}.

(2) La correspondance $s = 2p - \mathrm{id}$, $p = \frac{s +
\mathrm{id}}2$ est une \omterm{def:b1:logic:inj}{bijection} entre endomorphismes, et par elle
\[
s^2 = 4p^2 - 4p + \mathrm{id} = \mathrm{id}
\iff 4p^2 = 4p \iff p^2 = p :
\]
les \omterm{def:b1:linmaps:projection}{symétries} correspondent exactement aux \omterm{def:b1:linmaps:projection}{projections}. Traduisons les
\omterm{def:b1:vspaces:subspace}{sous-espaces}~: $s(x) = x \iff p(x) = x$, donc $\ker(s - \mathrm{id}) =
\operatorname{im} p = F$~; et $s(x) = -x \iff 2p(x) = 0 \iff x
\in \ker p = G$, donc $\ker(s + \mathrm{id}) = G$. La \omterm{def:b1:vspaces:sum}{somme directe}
$E = F \oplus G$ du point (1) devient la décomposition annoncée
en vecteurs fixes et vecteurs retournés de $s$.
\end{proof}

\begin{example}[Une projection et sa symétrie, explicitement]
\label{ex:b1:linmaps:projexplicit}
Dans $\R^2$, projetons sur $F = \operatorname{Vect}(1,1)$ parallèlement à $G =
\operatorname{Vect}(0,1)$. Décomposons $(x, y) = a(1,1) + b(0,1)$~:
la première coordonnée donne $a = x$, la seconde $b = y - x$.
D'où
\[
p(x, y) = (x, x),
\qquad
s(x, y) = 2p(x,y) - (x,y) = (x,\ 2x - y).
\]
Vérifions l'algèbre~: $p(p(x,y)) = p(x,x) = (x,x)$, et $s(s(x,y)) =
s(x, 2x - y) = (x, 2x - (2x - y)) = (x, y)$. Géométriquement, $s$
est la «~réflexion oblique~» par rapport à la droite $y = x$ dans la
direction verticale~: elle fixe $F$ point par point et retourne $G$. Si
nous avions projeté sur le même $F$ parallèlement à $G' =
\operatorname{Vect}(1,-1)$, la formule deviendrait
$p'(x,y) = \bigl(\frac{x+y}2, \frac{x+y}2\bigr)$~: une \omterm{def:b1:linmaps:projection}{projection} est
déterminée par son \omterm{def:b1:linmaps:kerim}{image} \emph{et} par son \omterm{def:b1:linmaps:kerim}{noyau}, jamais par la seule
\omterm{def:b1:linmaps:kerim}{image}.
\end{example}

\begin{omfigure}
\begin{tikzpicture}[scale=1.15]
  \draw[->, thin] (-0.4,0) -- (3.4,0) node[below] {$x$};
  \draw[->, thin] (0,-0.4) -- (0,3.8) node[left] {$y$};
  \draw[omDef, very thick] (-0.3,-0.3) -- (3.1,3.1)
    node[right, pos=0.98, font=\small] {$F:\ y = x$};
  \draw[omProp, thick, ->] (2.6,0.2) -- (2.6,0.9)
    node[right, font=\small] {$G$};
  \fill (2,0.5) circle (0.05) node[right, font=\small] {$M$};
  \fill (2,2) circle (0.05) node[left, font=\small] {$p(M)$};
  \fill (2,3.5) circle (0.05) node[left, font=\small] {$s(M)$};
  \draw[omThm, dashed] (2,0.5) -- (2,3.5);
\end{tikzpicture}

{\small La \omterm{def:b1:linmaps:projection}{projection} sur $F = \operatorname{Vect}(1,1)$ parallèlement à
$G = \operatorname{Vect}(0,1)$ et sa \omterm{def:b1:linmaps:projection}{symétrie}, sur le point $M
= (2,\ 0.5)$~: en glissant verticalement, $M$ atteint $F$ en $p(M) = (2,2)$
et arrive en $s(M) = 2p(M) - M = (2,\ 3.5)$, aussi loin au-dessus de $F$
(mesuré le long de $G$) que $M$ était en dessous.}
\end{omfigure}

\section{Formes linéaires et hyperplans}

\begin{definition}\label{def:b1:linmaps:forms}
Une \emph{forme linéaire}\index{forme linéaire} sur $E$ est une \omterm{def:b1:linmaps:def}{application
linéaire} $\varphi \colon E \to K$. Un \emph{hyperplan}\index{hyperplan} de
$E$ ($\dim E = n$) est un \omterm{def:b1:vspaces:subspace}{sous-espace} de dimension $n - 1$.
\end{definition}

\begin{example}[Une forme d'évaluation et son hyperplan]
\label{ex:b1:linmaps:evaluationform}
Sur $\R_2[X]$, l'évaluation $\varphi(P) = P(2)$ est une \omterm{def:b1:linmaps:forms}{forme
linéaire}, non nulle ($\varphi(1) = 1$). Son \omterm{def:b1:linmaps:kerim}{noyau} est l'\omterm{def:b1:linmaps:forms}{hyperplan}
des \omterm{def:b1:poly:def}{polynômes} s'annulant en $2$, c'est-à-dire (théorème de factorisation,
\cref{thm:b1:poly:factor}) les multiples de $X - 2$ à l'\omterm{def:b1:topology:closure}{intérieur} de
$\R_2[X]$~:
\[
\ker\varphi = \operatorname{Vect}\bigl(X - 2,\ X(X - 2)\bigr),
\qquad \dim = 2 .
\]
En \omterm{prop:b1:vspaces:coordinates}{coordonnées} dans la \omterm{def:b1:vspaces:free}{base} $(1, X, X^2)$, $\varphi(a + bX + cX^2) = a +
2b + 4c$~: toute \omterm{def:b1:linmaps:forms}{forme linéaire} sur un espace \omterm{def:b1:findim:def}{de dimension finie} est,
une fois une \omterm{def:b1:vspaces:free}{base} fixée, une expression \omterm{def:b1:linmaps:def}{linéaire} fixe en les
\omterm{prop:b1:vspaces:coordinates}{coordonnées} --- les formes sont des «~vecteurs lignes~», ce que le
\cref{ch:b1:matrices} rendra littéral, et la ligne de coefficients
obtenue ici, $(1, 2, 4)$, est une ligne de Vandermonde~: les formes
d'évaluation sont la porte par laquelle la théorie de l'interpolation du
devoir maison du \cref{ch:b1:det} entre en algèbre linéaire.
\end{example}

\begin{theorem}\label{thm:b1:linmaps:hyperplanes}
Les \omterm{def:b1:linmaps:forms}{hyperplans} de $E$ sont exactement les noyaux des \omterm{def:b1:linmaps:forms}{formes linéaires}
non nulles. Deux formes non nulles ont le même \omterm{def:b1:linmaps:kerim}{noyau} si et seulement si elles
sont proportionnelles.
\end{theorem}

\begin{proof}
Si $\varphi \neq 0$~: $\operatorname{rk}\varphi = 1$ (l'\omterm{def:b1:linmaps:kerim}{image} est un
\omterm{def:b1:vspaces:subspace}{sous-espace} non nul de $K$), donc $\dim\ker\varphi = n - 1$~: c'est un \omterm{def:b1:linmaps:forms}{hyperplan}.
Réciproquement, soient $H$ un \omterm{def:b1:linmaps:forms}{hyperplan}, $(e_1, \dots, e_{n-1})$ une \omterm{def:b1:vspaces:free}{base}
de $H$ complétée par $e_n$~: la forme «~dernière coordonnée~» a pour \omterm{def:b1:linmaps:kerim}{noyau}
$H$.

Deux formes proportionnelles ont le même \omterm{def:b1:linmaps:kerim}{noyau}. Réciproquement, supposons
$\ker\varphi = \ker\psi = H$ et choisissons $a \notin H$~: tout $x$ s'écrit
$x = h + \lambda a$ (car $E = H \oplus Ka$), et
\[
\varphi(x) = \lambda \varphi(a), \qquad \psi(x) = \lambda\psi(a):
\]
donc $\varphi = \frac{\varphi(a)}{\psi(a)}\,\psi$.
\end{proof}

\begin{example}\label{ex:b1:linmaps:hyperplane}
Dans $K^n$, un \omterm{def:b1:linmaps:forms}{hyperplan} est un \omterm{def:b1:logic:sets}{ensemble} de solutions $\{a_1 x_1 + \dots + a_n
x_n = 0\}$ avec les $a_i$ non tous nuls --- l'équation familière d'un
plan passant par l'origine dans $\R^3$. Dans les espaces de fonctions, les formes
d'évaluation $P \mapsto P(1)$ ou $f \mapsto \int_0^1 f$ définissent des \omterm{def:b1:linmaps:forms}{hyperplans}
de $\R_n[X]$, de $C(\intcc{0}{1})$ (cf.\ \cref{exo:b1:findim:6}).
\end{example}

\begin{example}[Un hyperplan, traité de trois façons]
\label{ex:b1:linmaps:hyperplanerun}
Prenons $\varphi(x, y, z) = x - 2y + 3z$ sur $\R^3$ et $H =
\ker\varphi$. \emph{\omterm{def:b1:vspaces:free}{Base}}~: résolvons $x = 2y - 3z$~:
\[
(2y - 3z,\ y,\ z) = y\,(2, 1, 0) + z\,(-3, 0, 1),
\]
deux vecteurs \omterm{def:b1:vspaces:free}{libres}~: $\dim H = 2$, un \omterm{def:b1:linmaps:forms}{hyperplan}, comme le
\cref{thm:b1:linmaps:hyperplanes} le prédit à partir de $\varphi \neq
0$. \emph{Droite \omterm{def:b1:vspaces:sum}{supplémentaire}}~: tout vecteur hors de $H$ en engendre une,
par exemple $a = (1, 0, 0)$ ($\varphi(a) = 1 \neq 0$)~; la
décomposition d'un $v$ arbitraire est explicite~:
\[
v = \underbrace{\bigl(v - \varphi(v)\,a\bigr)}_{\in\,H}
+ \underbrace{\varphi(v)\,a}_{\in\,\operatorname{Vect}(a)},
\]
puisque $\varphi\bigl(v - \varphi(v)a\bigr) = \varphi(v) -
\varphi(v)\varphi(a) = 0$. \emph{Proportionnalité}~: si $\psi(x,y,z)
= -2x + 4y - 6z$, alors $\psi = -2\varphi$ et toutes deux ont pour \omterm{def:b1:linmaps:kerim}{noyau}
$H$~; réciproquement, toute forme s'annulant sur $H$ est un multiple de
$\varphi$ (\cref{exo:b1:linmaps:8}) --- l'équation d'un
\omterm{def:b1:linmaps:forms}{hyperplan} est unique à un scalaire près, fait utilisé constamment pour les
plans en géométrie.
\end{example}

\begin{remark}[Pièges classiques]
\label{rem:b1:linmaps:pitfalls}
\emph{Le \omterm{def:b1:linmaps:kerim}{noyau} et l'\omterm{def:b1:linmaps:kerim}{image} vivent dans des espaces différents}~: $\ker u
\subseteq E$, $\operatorname{im} u \subseteq F$~; la somme $\ker u
+ \operatorname{im} u$ n'a de sens que pour les endomorphismes, et
même alors elle n'est pas nécessairement directe ($u(x, y) = (y, 0)$ vérifie $\ker u
= \operatorname{im} u$~; l'\cref{exo:b1:linmaps:7} caractérise
les cas où elle l'est).
\emph{$u^2 = 0$ ne signifie pas $u = 0$}~: le même $u(x,y) =
(y, 0)$ a un carré nul sans être nul --- ce que $u^2 = 0$
dit vraiment, c'est $\operatorname{im} u \subseteq \ker u$
(\cref{exo:b1:linmaps:5}).
\emph{L'équivalence \omterm{def:b1:logic:inj}{injective} $\iff$ \omterm{def:b1:logic:inj}{surjective} exige \omterm{def:b1:findim:def}{des dimensions finies}
égales}~: sur $K[X]$, la dérivation est \omterm{def:b1:logic:inj}{surjective} et non
\omterm{def:b1:logic:inj}{injective}, $P \mapsto XP$ \omterm{def:b1:logic:inj}{injective} et non \omterm{def:b1:logic:inj}{surjective}
(\cref{cor:b1:linmaps:samedim})~; et entre espaces de dimensions
\emph{différentes}, l'une des implications est tout simplement
impossible ($\operatorname{rk} u \leq \min(\dim E, \dim F)$).
\emph{Prescrire des \omterm{def:b1:linmaps:kerim}{images} marche sur une \omterm{def:b1:vspaces:free}{base}, pas sur une famille quelconque}~:
exiger $u(1, 0) = a$, $u(0, 1) = b$, $u(1, 1) = c$
surdétermine $u$ sauf si $c = a + b$~; une \omterm{def:b1:linmaps:def}{application linéaire} est \omterm{def:b1:vspaces:free}{libre} sur une
\omterm{def:b1:vspaces:free}{base}, asservie partout ailleurs.
\emph{Le \omterm{def:b1:linmaps:rank}{rang} n'est pas préservé par composition}~: il ne peut que chuter,
$\operatorname{rk}(vu) \leq \min(\operatorname{rk} u,
\operatorname{rk} v)$ (\cref{exo:b1:linmaps:4}), la perte exacte étant
mesurée dans l'\cref{exo:b1:linmaps:12}.
\end{remark}

\begin{remark}[Où ces applications mènent]
\label{rem:b1:linmaps:whereused}
Les \omterm{def:b1:linmaps:def}{applications linéaires} sont sur le point de devenir des \emph{matrices}~: une fois
des \omterm{def:b1:vspaces:free}{bases} fixées, le \cref{ch:b1:matrices} code tout $u \in \mathcal{L}(E,
F)$ par un tableau rectangulaire, et la composition devient le produit
matriciel --- le \omterm{thm:b1:linmaps:ranknullity}{théorème du rang} gouverne alors la théorie des systèmes
\omterm{def:b1:linmaps:def}{linéaires} du \cref{ch:b1:det}. Les \omterm{def:b1:linmaps:projection}{projections} reviennent au
\cref{ch:b1:euclid} dans leur cas particulier le plus utile, la \omterm{def:b1:linmaps:projection}{projection}
\emph{orthogonale}, où le \omterm{def:b1:linmaps:kerim}{noyau} est choisi perpendiculaire à l'\omterm{def:b1:linmaps:kerim}{image}.
Le devoir maison ci-dessous pousse l'algèbre des projecteurs aussi loin que
les outils de première année le permettent, jusqu'au \omterm{pb:b1:linmaps:1}{lemme de Fitting}~; le
volume de Licence~2 va plus loin avec la trace et la théorie des valeurs
propres, pour laquelle les projecteurs sur des \omterm{def:b1:vspaces:subspace}{sous-espaces} stables sont les
briques de \omterm{def:b1:vspaces:free}{base}.
\end{remark}

\begin{remark}[Perspectives à l'intérieur du livre 3~: le théorème du rang trois fois de plus]
\label{rem:b1:linmaps:perspectives}
La loi de conservation $\dim E = \dim\ker u + \operatorname{rk} u$
sera relue trois fois avant la fin du volume. Au
\cref{ch:b1:det} elle devient la forme des \omterm{def:b1:logic:sets}{ensembles} de solutions~: un
système compatible à $p$ inconnues et de \omterm{def:b1:linmaps:rank}{rang} $r$ a un \omterm{def:b1:logic:sets}{ensemble} de
solutions de dimension $p - r$ --- une dimension de \omterm{def:b1:linmaps:kerim}{noyau} déguisée. Au
\cref{ch:b1:euclid} elle se scinde orthogonalement, $\dim F + \dim
F^\perp = \dim E$, et alimente tout calcul de distance. Dans
le devoir maison du \cref{ch:b1:multivar}, elle est la
comptable des moindres carrés~: $n$ observations, $2$ paramètres ajustés,
$n - 2$ dimensions de résidus, et l'identité de Pythagore
$\norm b^2 = \norm p^2 + \norm{b - p}^2$ est l'ombre euclidienne du
\omterm{thm:b1:linmaps:ranknullity}{théorème du rang}. Un théorème, quatre costumes.
\end{remark}

\section{Exercices}

\begin{exercise}[$\star$]\label{exo:b1:linmaps:1}
Quelles \omterm{def:b1:logic:map}{applications} sont \omterm{def:b1:linmaps:def}{linéaires}~?
\begin{enumerate}
  \item $\R^2 \to \R^2$, $(x, y) \mapsto (x + y, x - 2y)$~;
  \item $\R^2 \to \R$, $(x, y) \mapsto xy$~;
  \item $\R[X] \to \R[X]$, $P \mapsto P' + XP$~;
  \item $\mathcal{F}(\R,\R) \to \R$, $f \mapsto f(3)$.
\end{enumerate}
\end{exercise}

\begin{exercise}[$\star$]\label{exo:b1:linmaps:2}
Soit $u \colon \R^3 \to \R^3$, $(x,y,z) \mapsto (x + y - z,\; 2x + y
+ z,\; 3x + 2y)$. Déterminer $\ker u$ (\omterm{def:b1:vspaces:free}{base}, dimension),
$\operatorname{rk} u$, et une \omterm{def:b1:vspaces:free}{base} de $\operatorname{im} u$. L'\omterm{def:b1:logic:map}{application} $u$
est-elle \omterm{def:b1:logic:inj}{injective}~? \omterm{def:b1:logic:inj}{surjective}~?
\end{exercise}

\begin{exercise}[$\star$]\label{exo:b1:linmaps:3}
Soit $u \colon \R_n[X] \to \R_n[X]$, $P \mapsto P - P'$. Démontrer que
$u$ est un isomorphisme~: une fois via $\ker u$, une fois en exhibant la
réciproque \emph{(considérer $P + P' + P'' + \dots$)}.
\end{exercise}

\begin{exercise}[$\star$]\label{exo:b1:linmaps:4}
Soient $u \in \mathcal{L}(E, F)$ et $v \in \mathcal{L}(F, G)$, les espaces
étant \omterm{def:b1:findim:def}{de dimension finie}. Démontrer~:
\[
\operatorname{rk}(v \circ u) \leq
\min\bigl(\operatorname{rk} u,\ \operatorname{rk} v\bigr).
\]
\end{exercise}

\begin{exercise}[$\star\star$]\label{exo:b1:linmaps:5}
Soit $u$ un endomorphisme de $E$ (\omterm{def:b1:findim:def}{de dimension finie}) avec $u^2 =
0$. Démontrer que $\operatorname{im} u \subseteq \ker u$, d'où
$\operatorname{rk} u \leq \frac{\dim E}{2}$. Pour $E = \R^2$, donner un
exemple où il y a égalité.
\end{exercise}

\begin{exercise}[$\star\star$]\label{exo:b1:linmaps:6}
Soient $p, q$ des \omterm{def:b1:linmaps:projection}{projections} de $E$ avec $p \circ q = q \circ p$.
Démontrer que $p \circ q$ est une \omterm{def:b1:linmaps:projection}{projection}, avec
\[
\operatorname{im}(pq) = \operatorname{im} p \cap \operatorname{im}
q ,
\qquad
\ker (pq) = \ker p + \ker q .
\]
\end{exercise}

\begin{exercise}[$\star\star$]\label{exo:b1:linmaps:7}
Soit $u \in \mathcal{L}(E)$, $E$ \omterm{def:b1:findim:def}{de dimension finie}. Démontrer
l'équivalence de~:
\begin{enumerate}
  \item $E = \ker u \oplus \operatorname{im} u$~;
  \item $\ker u = \ker u^2$~;
  \item $\operatorname{im} u = \operatorname{im} u^2$.
\end{enumerate}
\end{exercise}

\begin{exercise}[$\star\star$]\label{exo:b1:linmaps:8}
Soient $\varphi, \psi$ des \omterm{def:b1:linmaps:forms}{formes linéaires} sur $E$ avec $\ker\varphi
\subseteq \ker\psi$. Démontrer que $\psi = \lambda\varphi$ pour un certain
$\lambda \in K$ (y compris dans les cas dégénérés).
\end{exercise}

\begin{exercise}[$\star\star\star$]\label{exo:b1:linmaps:9}
Soit $u \in \mathcal{L}(E)$ avec $\dim E = n$, et supposons $u^n = 0$
mais $u^{n-1} \neq 0$ (un endomorphisme \emph{nilpotent maximal}).
Choisissons $x$ tel que $u^{n-1}(x) \neq 0$~; démontrer que $\bigl(x, u(x), \dots,
u^{n-1}(x)\bigr)$ est une \omterm{def:b1:vspaces:free}{base} de $E$. \emph{(Appliquer des puissances de $u$ à une
combinaison nulle, en commençant par $u^{n-1}$.)}
\end{exercise}

\begin{exercise}[$\star\star\star$]\label{exo:b1:linmaps:10}
Soit $f \in \mathcal{L}(\R^n)$ avec $f \circ f = -\mathrm{id}$.
\begin{enumerate}
  \item Démontrer que $f$ est un isomorphisme et qu'aucun $x \neq 0$
        ne vérifie $f(x) = \lambda x$ avec $\lambda \in \R$.
  \item Démontrer que $n$ est pair. \emph{Indication~: choisir $x_1 \neq 0$~; montrer
        que $\operatorname{Vect}(x_1, f(x_1))$ est un plan stable par
        $f$~; choisir $x_2$ hors de ce plan et itérer, en démontrant que
        $\bigl(x_1, f(x_1), x_2, f(x_2), \dots\bigr)$ reste \omterm{def:b1:vspaces:free}{libre}.}
\end{enumerate}
\end{exercise}

\begin{exercise}[$\star\star$]\label{exo:b1:linmaps:11}
Soient $u, v \in \mathcal{L}(E, F)$, les espaces étant \omterm{def:b1:findim:def}{de dimension finie}.
Démontrer l'encadrement
\[
\abs{\operatorname{rk} u - \operatorname{rk} v}
\;\leq\; \operatorname{rk}(u + v)
\;\leq\; \operatorname{rk} u + \operatorname{rk} v .
\]
\emph{(Pour la majoration, comparer $\operatorname{im}(u+v)$ avec
$\operatorname{im} u + \operatorname{im} v$~; pour la minoration,
appliquer astucieusement la majoration.)}
\end{exercise}

\begin{exercise}[$\star\star\star$]\label{exo:b1:linmaps:12}
(Inégalité de Frobenius) Soient $u \in \mathcal{L}(E, F)$, $w \in
\mathcal{L}(F, G)$ et $v \in \mathcal{L}(G, H)$, tous les espaces étant
\omterm{def:b1:findim:def}{de dimension finie}. Démontrer la formule exacte
\[
\operatorname{rk}(v \circ w) = \operatorname{rk} w -
\dim\bigl(\ker v \cap \operatorname{im} w\bigr),
\]
et en déduire l'inégalité de Frobenius
\[
\operatorname{rk}(v \circ w) + \operatorname{rk}(w \circ u)
\;\leq\; \operatorname{rk} w + \operatorname{rk}(v \circ w \circ
u) .
\]
Vérifier que le cas $w = \mathrm{id}_F$ est l'inégalité de Sylvester,
démontrée sous forme matricielle dans l'\cref{exo:b1:matrices:10}.
\end{exercise}

\section{Problème~: calcul sur les projecteurs et lemme de Fitting}

\begin{problem}\label{pb:b1:linmaps:1}
Les \omterm{def:b1:linmaps:projection}{projections} sont les endomorphismes que les \omterm{def:b1:vspaces:sum}{sommes directes} produisent, et
réciproquement~: toute identité $E = F_1 \oplus \dots \oplus F_k$ est
secrètement une famille de projecteurs de somme l'identité. Ce
problème développe ce dictionnaire --- l'algèbre d'un
projecteur, de deux, de $k$ --- puis applique les mêmes idées de
stabilisation à un endomorphisme quelconque et démontre le
\emph{\omterm{pb:b1:linmaps:1}{lemme de Fitting}}~: tout endomorphisme d'un espace \omterm{def:b1:findim:def}{de
dimension finie} se scinde en une partie nilpotente et une
partie inversible. Dans tout ce qui suit, $E$ est un $K$-espace vectoriel de
dimension $n$, et \emph{projecteur} signifie $p \in \mathcal{L}(E)$
avec $p^2 = p$ (\cref{thm:b1:linmaps:projchar}).
\index{projection}\index{lemme de Fitting}\index{nilpotent}

\medskip
\noindent\textbf{Partie I --- L'algèbre autour d'un projecteur.}
Soit $p$ un projecteur, $p \neq 0$, $p \neq \mathrm{id}$.
\begin{enumerate}
  \item Montrer que $\mathrm{id} - p$ est un projecteur et identifier
        $\operatorname{im}(\mathrm{id} - p)$ et
        $\ker(\mathrm{id} - p)$.
  \item Calculer $(\lambda\,\mathrm{id} + \mu\,p)^2$ et déterminer
        tous les couples $(\lambda, \mu) \in K^2$ pour lesquels
        $\lambda\,\mathrm{id} + \mu\,p$ est un projecteur.
  \item Montrer que le plan $\operatorname{Vect}(\mathrm{id}, p)$
        de $\mathcal{L}(E)$ est stable par composition, et que
        pour tout \omterm{def:b1:poly:def}{polynôme} $Q \in K[X]$,
        \[
        Q(p) = Q(0)\,\mathrm{id} + \bigl(Q(1) -
        Q(0)\bigr)\,p .
        \]
  \item Déterminer pour quels $(\lambda, \mu)$ l'\omterm{def:b1:logic:map}{application}
        $\lambda\,\mathrm{id} + \mu\,p$ est inversible, et donner
        son inverse sous la forme $\alpha\,\mathrm{id} + \beta\,p$.
        Interpréter la réponse par l'action de
        $\lambda\,\mathrm{id} + \mu\,p$ sur $\operatorname{im} p$
        et sur $\ker p$.
  \item Soit $p'$ un autre projecteur ayant la \emph{même \omterm{def:b1:linmaps:kerim}{image}}
        $\operatorname{im} p' = \operatorname{im} p$. Montrer que
        $p\,p' = p'$ et $p'\,p = p$. Que disent ces identités
        sur la composition de \omterm{def:b1:linmaps:projection}{projections} sur le même \omterm{def:b1:vspaces:subspace}{sous-espace}
        parallèlement à des noyaux différents~?
\end{enumerate}

\medskip
\noindent\textbf{Partie II --- Deux projecteurs.} Soient $p, q$ des
projecteurs de $E$~; on suppose la caractéristique différente de $2$ (vrai
pour $K = \R, \C$).
\begin{enumerate}[resume]
  \item Supposons que $p + q$ soit un projecteur. En développant $(p + q)^2$,
        montrer que $pq + qp = 0$~; en composant avec $p$ à gauche, puis
        à droite, en déduire que $pq = qp$, et conclure que $pq = qp =
        0$.
  \item Réciproquement, supposons $pq = qp = 0$. Montrer que $p + q$ est un
        projecteur, avec
        \[
        \operatorname{im}(p + q) = \operatorname{im} p \oplus
        \operatorname{im} q,
        \qquad
        \ker(p + q) = \ker p \cap \ker q .
        \]
  \item Montrer que $p - q$ est un projecteur si et seulement si $pq = qp
        = q$. \emph{(Appliquer les questions 6 et 7 à $\mathrm{id} - p$
        et $q$.)}
  \item Montrer la signification géométrique de $pq = qp = q$~: elle a lieu si
        et seulement si $\operatorname{im} q \subseteq
        \operatorname{im} p$ et $\ker p \subseteq \ker q$. (On
        écrit alors $q \leq p$~: «~$q$ projette sur moins, parallèlement à
        plus~».)
  \item Supposons maintenant que $p$ et $q$ commutent. Rappelons, d'après
        l'\cref{exo:b1:linmaps:6}, que $pq$ est le projecteur sur
        $\operatorname{im} p \cap \operatorname{im} q$ parallèlement à
        $\ker p + \ker q$. Montrer que $r = p + q - pq$ est un
        projecteur avec
        \[
        \operatorname{im} r = \operatorname{im} p +
        \operatorname{im} q,
        \qquad
        \ker r = \ker p \cap \ker q .
        \]
        \emph{(Considérer $\mathrm{id} - r = (\mathrm{id} -
        p)(\mathrm{id} - q)$.)}
\end{enumerate}

\medskip
\noindent\textbf{Partie III --- Décompositions de l'identité.}
\begin{enumerate}[resume]
  \item Soit $E = F_1 \oplus \dots \oplus F_k$ et, pour $x = x_1 +
        \dots + x_k$ (décomposition unique, $x_i \in F_i$), posons
        $p_i(x) = x_i$. Montrer que chaque $p_i$ est un projecteur, que
        $p_i p_j = 0$ pour $i \neq j$, et que $p_1 + \dots + p_k
        = \mathrm{id}$~; identifier $\operatorname{im} p_i$ et
        $\ker p_i$.
  \item Réciproquement, soient $p_1, \dots, p_k \in \mathcal{L}(E)$
        vérifiant $p_1 + \dots + p_k = \mathrm{id}$ et $p_i p_j =
        0$ pour tous $i \neq j$. Montrer que chaque $p_i$ est un
        projecteur et que $E = \operatorname{im} p_1 \oplus
        \dots \oplus \operatorname{im} p_k$.
  \item Deux projecteurs avec $p + q = \mathrm{id}$~: montrer que $pq
        = qp = 0$ est automatique.
  \item Trois projecteurs avec $p + q + r = \mathrm{id}$~: montrer
        que $p + q$ est un projecteur, et déduire de la question 6
        que \emph{tous} les produits deux à deux sont nuls --- d'où $E =
        \operatorname{im} p \oplus \operatorname{im} q \oplus
        \operatorname{im} r$, sans aucune hypothèse sur les produits.
  \item Pour $k$ projecteurs avec $p_1 + \dots + p_k =
        \mathrm{id}$~: montrer d'abord que, pour des \omterm{def:b1:vspaces:subspace}{sous-espaces} quelconques, $\dim
        (F_1 + \dots + F_k) \leq \dim F_1 + \dots + \dim F_k$,
        avec égalité si et seulement si la somme est directe~; montrer
        ensuite que $E = \operatorname{im} p_1 + \dots +
        \operatorname{im} p_k$, et démontrer que \emph{si}
        de plus $\sum_i \operatorname{rk} p_i \leq n$, la somme
        est directe et $p_i p_j = 0$ pour $i \neq j$.
\end{enumerate}

\medskip
\noindent\textbf{Partie IV --- Noyaux itérés~: le \omterm{pb:b1:linmaps:1}{lemme de Fitting}.}
Soit $u \in \mathcal{L}(E)$, $\dim E = n$.
\begin{enumerate}[resume]
  \item Montrer les deux chaînes, valables pour tout $k \geq 0$~:
        \[
        \ker u^k \subseteq \ker u^{k+1},
        \qquad
        \operatorname{im} u^{k+1} \subseteq \operatorname{im}
        u^k .
        \]
  \item Montrer que si $\ker u^{r} = \ker u^{r+1}$ pour un certain $r$,
        alors $\ker u^{k} = \ker u^{r}$ pour tout $k \geq r$~; énoncer
        et démontrer la stabilisation analogue pour les \omterm{def:b1:linmaps:kerim}{images}.
  \item En déduire qu'il existe un plus petit entier $r$ tel que $\ker
        u^{r} = \ker u^{r+1}$, que $r \leq n$, et que les \omterm{def:b1:linmaps:kerim}{images}
        se stabilisent au même $r$.
  \item (\omterm{pb:b1:linmaps:1}{Lemme de Fitting}) Démontrer que
        \[
        E \;=\; \ker u^{r} \,\oplus\, \operatorname{im} u^{r} .
        \]
  \item Montrer que les deux \omterm{def:b1:vspaces:subspace}{sous-espaces} sont stables par $u$, que la
        restriction de $u$ à $\ker u^{r}$ est nilpotente, et que
        la restriction de $u$ à $\operatorname{im} u^{r}$ est un
        isomorphisme de $\operatorname{im} u^{r}$~: tout
        endomorphisme est, sur une \omterm{def:b1:vspaces:sum}{somme directe} canonique, «~nilpotent
        plus inversible~».
  \item Soit $\pi$ le projecteur sur $\ker u^{r}$ parallèlement à
        $\operatorname{im} u^{r}$. Montrer que $\pi \circ u = u
        \circ \pi$.
\end{enumerate}

\medskip
\noindent\textbf{Partie V --- Un cas traité, et synthèse.}
\begin{enumerate}[resume]
  \item Soit $u(x, y, z) = (y, 0, z)$ sur $\R^3$. Calculer $u^2$ et
        $u^3$, déterminer l'indice de stabilisation $r$, les
        \omterm{def:b1:vspaces:subspace}{sous-espaces} $\ker u^{r}$ et $\operatorname{im} u^{r}$, le
        projecteur de Fitting $\pi$, et vérifier sur les formules que
        $\pi u = u\pi$ et que $u$ est nilpotente sur un facteur,
        \omterm{def:b1:logic:inj}{bijective} sur l'autre.
  \item Montrer les équivalences~: $u$ nilpotente $\iff$ $\ker u^{r} =
        E$ $\iff$ $\pi = \mathrm{id}$~; et en déduire qu'un
        endomorphisme nilpotent d'un espace de dimension $n$
        vérifie toujours $u^{n} = 0$ (l'indice de nilpotence ne
        dépasse jamais la dimension).
  \item (Unicité) Supposons $E = A \oplus B$ avec $A, B$ stables
        par $u$, la restriction $u|_A$ nilpotente et $u|_B$
        \omterm{def:b1:logic:inj}{bijective}. Démontrer que $A = \ker u^{r}$ et $B =
        \operatorname{im} u^{r}$~: la décomposition de Fitting est
        unique.
  \item Synthèse, en quatre phrases~: quel dictionnaire la partie III
        établit entre \omterm{def:b1:vspaces:sum}{sommes directes} et familles de
        projecteurs~; pourquoi la question 14 n'a eu besoin d'aucune
        hypothèse sur les produits alors que la question 15 a eu besoin d'une hypothèse de \omterm{def:b1:linmaps:rank}{rang}
        (et quel outil de Licence~2, la trace, la supprime)~; en quel
        sens le \omterm{pb:b1:linmaps:1}{lemme de Fitting} est la version stabilisée de
        l'\cref{exo:b1:linmaps:7}~; et ce que deviennent les deux facteurs de
        Fitting dans la théorie des valeurs propres du volume de
        Licence~2. Nommer le théorème démontré dans la partie IV.
\end{enumerate}
\end{problem}
